\section{The two residual root witnesses in the tensor seed}

Fix a Case III tensor datum $d$ of the source type $\mathrm{Dat}$,
over the original algebraic closure $\Omega$. Its arithmetic inputs
include an initial tower $T_0$, a quadratic field $K'$, the integral
rank-two matrices $\rho_1,\tau$ over the enlarged coefficient ring
$O_1$, and a twisting character $\delta$ over $O\subset O_1$.
Write $\Phi=\chi_\ell^{k-1}\phi$ for the first factor's residual
diagonal character before reduction; here
$\phi=\xi\psi^{\pm1}$ with sign prescribed by the label of $T_0$.
On the group of each tower $T$, the specified pair $Q_T$ restricts
$\delta\,\mathrm{tn}(\rho_1,\tau)$, where $\mathrm{tn}$ is the signed
tensor matrix in the source basis. At a place $A$ over $\ell$, the
two integral factors supply a slope $x_A$ and a sign $s_A=\pm1$.
The flag of $Q_T$ at $v$ uses $A=T.\mathrm{pl}(v)$.

% BEGIN E65S1
\paragraph{E65S1. The four residual characters and every local type-$A$ flag.}
Take any $T\ge T_0$, auxiliary set $W$, field $k'$ and ring map
$f:O_1\to k'$ killing the maximal ideal. Bars denote passage to this
field, including the inclusion $O\to O_1$ as necessary. This inclusion
preserves the maximal ideals (both coefficient rings use the fixed
$\ell$-adic valuation), so equal residues over $O$ have equal images.
For $g\in G_{\mathbb Q}$ put $P(g)=\bar\Phi(g)$, $b(g)=\bar\rho_1(g)_{10}$
and let $e(g)=1$ on $G_{K'}$ and $-1$ on the other coset. The defining
congruences and the two integral coset formulas for $\tau$ give
\[
 \bar\rho_1(g)=\begin{pmatrix}P(g)&0\\ b(g)&1\end{pmatrix},
 \qquad \bar\tau(g)=\operatorname{diag}(1,e(g)).
\]
In particular $f(\det\rho_1(g))=P(g)$ on the entire absolute group.
In the source basis the residual tensor matrix and local flag are
\[
 \bar\delta(g)
 \begin{pmatrix}
 P&0&0&0\\0&Pe&0&0\\0&-be&e&0\\b&0&0&1
 \end{pmatrix},\qquad
 n(x,s)=\begin{pmatrix}
 1&0&0&0\\s&1&0&0\\-xs&-x&1&0\\x&0&-s&1
 \end{pmatrix},
\]
respectively; the second formula uses $x=\bar x_A$, $s=\bar s_A$.
These follow by substitution in the signed tensor and flag matrices.
On $G_{T.F}$ the label is that of $T_0$. The twist gives
$\bar\chi_0=\bar\delta P$ and $\bar\chi_3=\bar\delta$, and the
quadratic condition gives $\bar\chi_0/\bar\chi_1=e$ by restriction
from $G_{T_0.F}$. Since $\chi_0\chi_3=\chi_1\chi_2$ and $e^2=1$,
the middle diagonal values are precisely $\bar\chi_1,\bar\chi_2$.
Conjugation by the lower unitriangular $n(x,s)$ preserves the diagonal
of this lower triangular representation; hence these are the required
local characters as well. The only possible off-diagonal weights are
those of $21,30$, namely $-\beta,-\tau$ in the notation of UX.46.
They lie in both $N_M$ and the type-$A$ support, and every local $10$
flag entry is a sign, hence nonzero. This supplies the residual
diagonal, local character, and shape fields of $(T.\mathrm{fr}\,W).\mathrm{Res}$,
both support assertions, and the determinant identity in $\mathrm{ResOf4}$.
% END E65S1

% BEGIN E65S2
\paragraph{E65S2. One element witnesses both negative roots on each acceptable tower.}
Continue with $T\ge T_0$ and now also assume $T.\mathrm{Ok}$. Put
$h_0=\operatorname{res}_{\mathbb Q}(T.g_0)$. The marked eigenvalues
with indices $0,3$ differ modulo the maximal ideal by acceptability.
Their ratio is the residual value of $\Phi(h_0)$ by E65S1.
Consequently $P(h_0)\ne1$ for the given $f$, since the nonzero
difference in the original residue field maps injectively into $k'$.
The integral lattice hypothesis on $\rho_1$ supplies an absolute $g$
with $(\rho_1(g))_{10}$ a unit, so $b(g)\ne0$.

The lower left entry in E65S1 satisfies
$b(uv)=b(u)P(v)+b(v)$. If $b(h_0)\ne0$, take $h=h_0$. Otherwise
put $h=gh_0g^{-1}$. Computing $b(hg)=b(gh_0)$ in the two orders gives
\[
           b(h)P(g)=b(g)(P(h_0)-1)\ne0 .
\]
The field $T.F$ is Galois over $\mathbb Q$, so $h$ still belongs to
$G_{T.F}$ under the original restriction map. At this one element the
$21$ and $30$ entries are $-\bar\delta(h)b(h)e(h)$ and
$\bar\delta(h)b(h)$, respectively. Both are nonzero, as $e(h)$ and
$\bar\delta(h)$ are units. This proves $\mathrm{EntOf4}$ on every
acceptable later tower and for every permitted coefficient-field map.
% END E65S2

% BEGIN E65S3
\paragraph{E65S3. The chosen local flags descend along the original tower map.}
Take $T\ge T_0$ and $f$ as in E65S1. The representations are
restrictions of the same rational one. For a place $v'$ of $T$ above
$\ell$, the covering property of $T_0$ gives a place $v$ of $T_0$
and $g\in G_{T_0.F}$ moving $A=T_0.\mathrm{pl}(v)$ to
$A'=T.\mathrm{pl}(v')$. Identify $g$ with its rational image.
We prove that $\bar Q_{T_0}.r(g)$ carries the source flag at $A$
to the flag at $A'$ modulo an upper triangular matrix.

There is an element $h\in D_{A'}$ with $P(h)\ne1$. Indeed the
original Frobenius lift $D.\mathrm{sg}(\ell)$ at $D.\mathrm{pl}(\ell)$
has cyclotomic value one. There $\Phi$ has the same value as
$\xi\Psi^{\pm1}$, whose reduction is nontrivial on that element by
the goodness condition $D.\mathrm{condD}$. Conjugating to $A'$
preserves the value of the character, and its nontriviality survives
$f$. Put $h'=g^{-1}hg\in D_A$. The rank-two local flag equations
make $(1,\bar x_A)^t$ a $P(h')$-eigenvector for $\bar\rho_1(h')$
and $(1,\bar x_{A'})^t$ a $P(h)$-eigenvector for $\bar\rho_1(h)$.
The two eigenvalues at $h$ are $P(h)\ne1$ and $1$. Intertwining
by $\bar\rho_1(g)$ and comparing first coordinates thus gives
\[
             b(g)+\bar x_A=P(g)\bar x_{A'} .
\]
The sign transport in the second factor gives
$\bar s_{A'}=e(g)\bar s_A$.

For the matrix check let $L(z)=\left(\begin{smallmatrix}1&0\\z&1\end{smallmatrix}\right)$
and $N(s)=1+s(E_{10}-E_{32})$. The displayed source flag in E65S1
equals $\mathrm{tn}(L(x),1)N(s)$. Our slope identity says that
$L(-\bar x_{A'})\bar\rho_1(g)L(\bar x_A)=\operatorname{diag}(P(g),1)$.
By multiplicativity of the tensor matrix, conjugation between the two
flags therefore gives, up to the unit twist,
\[
 N(-\bar s_{A'})\operatorname{diag}(P,Pe,e,1)N(\bar s_A).
\]
The only possibly nonzero strict lower entries are $10$ and $32$.
They are $P(e\bar s_A-\bar s_{A'})$ and
$e\bar s_{A'}-\bar s_A$, both zero. This proves the upper triangular
comparison for each actual place mover chosen above and hence the full
$\mathrm{OverOf4}$ compatibility.
% END E65S3

% BEGIN E65S4
\paragraph{E65S4. Continuity before and after specialization.}
This part holds on any tower $T$ for the original datum. The
entries of $\rho_1,\tau$ are locally constant modulo every power of
the maximal ideal of $O_1$. The same is true of $\delta$ after the
coefficient inclusion, since it takes the $n$th power of the maximal
ideal of $O$ into the $n$th power for $O_1$. Each entry of the
twisted tensor is a signed product of one entry from each factor
and a value of $\delta$. Modulo any prescribed power, that product
is locally constant. Restriction to $G_{T.F}$ is continuous, proving
the source adic continuity of $Q_T$. Finally, for any map $f:O_1\to R$
to a commutative ring killing the maximal ideal, each specialized
entry factors through its locally constant reduction modulo that
ideal. It is therefore locally constant with the discrete target,
giving continuity at the zero ideal as required by
$\mathrm{ContOf4}$. These four statements verify properties of the
tensor seed attached to the supplied arithmetic datum $d$.
% END E65S4
